\chapter{Estado del arte}
\label{cap:estadoDelArte}

Este capítulo revisa los trabajos existentes más relevantes en relación
con los objetivos formulados en el Capítulo~\ref{cap:introduccion}. La
revisión se organiza del siguiente modo. En primer lugar se revisa el
modelado de pricings de SaaS en el canal web, de donde proceden el
metamodelo y las herramientas cuyo patrón replica después el canal API. A
continuación, el modelado de la capacidad de las APIs y el análisis de esa
capacidad. Después, la extensión del ecosistema de Pricing Intelligence al
canal API: el ACM, la Datasheet y su motor de análisis, que son la
plataforma inmediata sobre la que este trabajo se construye. Por último,
el dominio de aplicación, las APIs de correo electrónico.
El capítulo se cierra con un
análisis del estado del arte que sintetiza qué resuelve cada propuesta y
qué limitaciones concretas permanecen abiertas, de las que se derivan
las contribuciones enumeradas en la
Sección~\ref{sec:introduccion:contribuciones}.

Conviene advertir una particularidad de esta revisión. El bloque
dedicado al canal API procede de dos Trabajos Fin de Máster coordinados
del mismo grupo de investigación. Uno aporta el modelo de capacidad, su
serialización y el motor de análisis, y se defendió pocos meses antes
que este trabajo \citep{gavira2026tfm}. El otro, desarrollado en
paralelo, se centra en el agente conversacional del ecosistema y
comparte la autoría del modelo \citep{ruizlopez2026tfm}. Ninguno de los
dos es aquí una referencia lejana. Ambos son el punto de partida
operativo de este trabajo, y por eso se revisan con un detalle mayor del
habitual. Una de las preguntas de investigación de este TFM, la de la
expresividad de ACM2YAML, responde a una línea que Gavira Sánchez deja
pendiente de forma explícita \citep{gavira2026tfm}. Las demás se derivan
del análisis de la Sección~\ref{sec:estado:analisis}.

\section{Modelado de pricings de SaaS}
\label{sec:estado:pricingsaas}

El pricing de un servicio SaaS es la estructura de planes,
características (\textit{features}), complementos (\textit{add-ons}) y
límites de uso con la que el proveedor regula el acceso a su servicio.
Ha dejado de ser un elemento puramente comercial para convertirse en un
artefacto que condiciona el diseño, la operación y la evolución del
propio sistema. La evidencia empírica de este desplazamiento la aportan
los estudios de García-Fernández et al. El primero analiza más de 150
pricings de una treintena de SaaS comerciales a lo largo de seis años y
documenta espacios de configuración de miles de combinaciones y un
crecimiento sostenido de su complejidad
\citep{garciafernandez2024racing}. El segundo es un estudio posterior
sobre el grado de soporte industrial para la operación dirigida por
pricings (\textit{pricing-driven DevOps}), publicado en \textit{IEEE
Transactions on Services Computing} \citep{garciafernandez2025trends}.
Gestionar manualmente estructuras de ese tamaño no es viable, y esa
constatación es la que motiva su formalización.

La respuesta del grupo ISA a este problema sigue una tradición que le es
propia: convertir un artefacto informal en un modelo estructurado para
someterlo a análisis automatizado. Es el mismo camino que en su día
recorrieron los modelos de características \citep{benavides2010fm}. Su
materialización para los pricings es el metamodelo \textbf{iPricing},
propuesto originalmente bajo el nombre Pricing4SaaS
\citep{garciafernandez2024pricing4saas} y consolidado después con su
denominación actual \citep{garciafernandez2025ipricing}. El metamodelo
estructura un pricing en planes, complementos y características. Estas últimas están
sujetas a límites de uso renovables o no renovables. Su serialización,
Pricing2YAML, permite representar pricings de forma portable e
interoperable \citep{garciafernandez2025trends}. Sobre esta base, el
trabajo presentado en CAiSE~2025 define siete operaciones de análisis
(cardinalidad del espacio de configuración, filtrado, enumeración de
suscripciones, coste, validación del pricing y de la suscripción, y
suscripción óptima). Las siete se resuelven traduciendo el pricing a un
problema de satisfacción de restricciones con optimización (CSOP). El
trabajo las valida sobre más de 150 pricings reales y detecta
inconsistencias en 35 de ellos \citep{garciafernandez2025ipricing}.

En torno a iPricing se ha construido un ecosistema de herramientas
interoperables, articulado por la línea de \textit{Pricing
Intelligence} \citep{cavero2026pricingintelligence}. A-MINT automatiza
la extracción: toma la URL de la página de pricing de un SaaS y devuelve
un iPricing válido en Pricing2YAML. Lo hace con un pipeline de modelos de
lenguaje que valida y corrige su propia salida en bucle. Alcanza una F1
del 70,51\,\% y valida con éxito el 98,73\,\% de los modelos que genera
\citep{cavero2025amint}, evaluado sobre SPECTRA, un corpus público con
los pricings de una treintena de SaaS entre 2019 y 2024
\citep{spectra2025}.
PRIME (\textit{PRicing Intelligence and Modelling Engine}) es el motor
que ejecuta las operaciones de análisis sobre el modelo
\citep{garciafernandez2025ipricing}. HARVEY, presentado en CAiSE~2026,
es el agente conversacional que traduce preguntas en lenguaje natural a
llamadas sobre esas herramientas y sintetiza los resultados
\citep{cavero2026pricingintelligence}. SPHERE, por último, actúa como
\textit{hub} del ecosistema: repositorio compartido de pricings
organizados en colecciones, editor con validación formal y renderizado
en tiempo real, y punto de acceso público a las demás herramientas
\citep{sphere}.

De este bloque interesan dos rasgos para los objetivos de este trabajo.
El primero, el papel del artefacto central. El pricing formalizado
concentra la información en una representación común e independiente de
la implementación. Sobre esa fuente compartida de conocimiento operan
extracción, análisis, consulta y publicación, y el canal API replicará
después ese patrón con la Datasheet. El segundo, la especialización de
canal. iPricing fue concebido para el canal web y no representa conceptos
del canal API como los \textit{endpoints}, los \textit{rates} por ventana
temporal o las reglas de consumo por petición. No es una deficiencia:
Gavira Sánchez concluye de forma explícita que canales distintos
requieren modelos distintos \citep{gavira2026tfm}. Aun así, esa
especialización delimita dónde termina su alcance y dónde debe empezar
otro instrumental.

\section{Modelado de la capacidad de las APIs}
\label{sec:estado:capacidad}

En el canal API la situación de partida es asimétrica. La superficie
funcional de una API REST \citep{fielding2000rest} cuenta con un
estándar de descripción maduro y universalmente adoptado, la
especificación OpenAPI \citep{openapi2025}, que documenta rutas,
métodos, parámetros y respuestas. Las condiciones de consumo, en cambio,
carecen de un equivalente asentado: cuántas peticiones admite un
endpoint por unidad de tiempo, qué volumen mensual permite cada
plan, qué ocurre al superarlo. Y es la información determinante para
quien integra el servicio.

Que esas condiciones son ubicuas y heterogéneas está documentado
empíricamente. El análisis de Gámez-Díaz et al. sobre 69 ofertas comerciales de APIs
RESTful muestra que las cuotas y los \textit{rates} son las limitaciones
dominantes y que la métrica más habitual es el número de peticiones
\citep{gamezdiaz2017offerings}. Muestra también que los modelos de
negocio se articulan en planes escalonados, con y sin sobrecoste por
exceso. Un estudio posterior, realizado junto a actores industriales del
sector de la gestión de APIs, sitúa esas limitaciones y los acuerdos de
nivel de servicio que las contienen en el centro de la industria del API
\citep{gamezdiaz2019limitations}. El mismo estudio constata la ausencia
de una representación estándar y procesable de esas limitaciones.

Los intentos de dotar a esas condiciones de una representación formal
heredan de la larga tradición de especificación de acuerdos de nivel de
servicio, cuyo exponente clásico es WS-Agreement \citep{andrieux2007wsagreement}.
SLA4OAI traslada esa línea al ecosistema
OpenAPI como extensión de la especificación \citep{gamezdiaz2019sla4oai}. Ancla las limitaciones
(cuotas y \textit{rates}, con umbrales, periodos y costes por operación
o por exceso) a las mismas rutas y métodos que el documento funcional ya
describe. El respaldo conceptual de esa sintaxis es
\textbf{Pricing4APIs}, el modelo riguroso de pricings de APIs RESTful de
Fresno-Aranda et al. \citep{fresno2025pricing4apis}. Organiza la
información en planes, limitaciones y capacidad, entendida como el techo
máximo que el proveedor puede ofrecer, del que las limitaciones son
fracciones. El modelo incorpora además una operación de validación con
criterios jerárquicos (del límite individual al pricing completo) que
detecta inconsistencias como límites inalcanzables o planes más caros
menos generosos que otros más baratos \citep{fresno2025pricing4apis}.

Tres de los rasgos que Gavira Sánchez destaca en esta familia de
propuestas \citep{gavira2026tfm} resultan determinantes para este
trabajo, porque son los que motivaron el modelo que se revisa en la
Sección~\ref{sec:estado:acm}. Primero, la capacidad es
unidimensional. Se reduce a un valor máximo de \textit{throughput} sobre
la métrica de peticiones, de modo que restricciones reales como la
concurrencia o el tamaño de la respuesta quedan fuera de su semántica.
Segundo, la distinción entre cuota y \textit{rate} se delega en el
\textit{limitador}, el componente del proveedor que lleva la cuenta de
las peticiones recibidas y rechaza las que exceden lo permitido. Que un
límite sea cuota o \textit{rate} depende de cómo ese componente cuente:
sobre una ventana de tiempo fija o sobre una ventana deslizante. Es, en
palabras de Gavira Sánchez, «un detalle de implementación que la mayoría
de proveedores no publica» \citep{gavira2026tfm}. Dos
límites publicados pueden resultar inclasificables desde la
documentación. Y tercero, el plan es el único eje de segmentación:
cualquier variación en las condiciones de dos consumidores exige planes
distintos \citep{fresno2025pricing4apis,gavira2026tfm}.

\section{Análisis de la capacidad de las APIs}
\label{sec:estado:analisis-capacidad}

Modelar la capacidad es condición necesaria pero no suficiente: el valor
está en poder razonar sobre ella. La gestión de la capacidad es una
preocupación clásica de la gestión de servicios TI
\citep{grummit2010capacity}. Trasladada a las APIs cristaliza en las
\textit{operaciones de análisis de la capacidad}, cómputos formales
sobre una noción de capacidad que evalúan el uso de una API bajo
condiciones de carga determinadas. En la literatura coexisten dos
aproximaciones, que difieren en cómo entienden la capacidad y en cómo
resuelven el análisis.

La primera, de Fresno-Aranda et al., opera sobre arquitecturas de
microservicios \citep{fowler2014microservices} que consumen APIs
externas reguladas por planes con limitaciones, bautizadas LAMA
(\textit{Limitation-Aware Microservice Architecture}). La capacidad es
el máximo volumen de trabajo que la arquitectura completa puede atender
sin exceder ninguna limitación externa. Smart LAMA, la API REST que los
autores construyen sobre ese estilo, traduce la topología (un grafo de servicios propios y APIs
consumidas) a un CSOP. Resolverlo responde a operaciones como el máximo
de peticiones servibles en un periodo y coste dados, el coste mínimo
para un volumen dado o el tiempo mínimo para completarlo
\citep{fresno2023gecon,fresno2022lama,fresno2025tesis}. La misma línea
aporta además un catálogo de once operaciones de análisis centradas en
el consumidor individual de una API \citep{fresno2023catalogo}. Sus
cinco escenarios van del suscriptor que calcula tiempos de espera al
prosumidor que busca la suscripción que soporte su propio modelo de
precios. Para Gavira Sánchez, frente a Smart LAMA, el catálogo convierte
por primera vez el pricing de una API en objeto directo del análisis,
aunque se definió sin implementación de referencia \citep{gavira2026tfm}.

La segunda aproximación es geométrica y parte de la visión de la capacidad
bajo un enfoque de caja negra planteada por Antonio Ruiz-Cortés y
Miguel Ruiz-López \citep{ruizcortes2025cajanegra}. El \textit{rate} se
interpreta como la velocidad máxima de consumo y la cuota como el
volumen máximo acumulable en un periodo. Ambos límites operan a la vez, y su intersección define la capacidad
efectiva.
Esta idea se materializa en la biblioteca Pricing4API, cuyo primer
prototipo desarrolló Daniel Ruiz López y que Gavira Sánchez extendió en su
Trabajo Fin de Grado \citep{pricing4api,gavira2025tfg}. El
\textit{BoundedRate} agrupa un rate y una o varias cuotas
interdependientes. Sobre él se define la \textbf{curva de capacidad}
$C(t)$: el máximo número de peticiones realizables desde el instante
inicial hasta $t$ sin violar ninguna restricción. Su ventaja es doble.
Se lee de un vistazo en una gráfica, y evaluarla en un punto cuesta mucho
menos que resolver un problema de restricciones. Sobre ella la biblioteca
responde preguntas de tiempo mínimo, capacidad acumulada, agotamiento de
cuota o rentabilidad entre planes \citep{gavira2025tfg}.

Leídas desde los objetivos de este TFM, ambas aproximaciones comparten
dos características que delimitan su alcance. La primera es la entrada.
Smart LAMA asume los pricings ya descritos conforme a Pricing4APIs, y la
biblioteca Pricing4API ni siquiera contaba con un artefacto de entrada
estructurado: cada plan debía codificarse a mano en Python, trasladando
uno a uno los parámetros de la documentación del proveedor. Ninguna de
las dos responde de dónde sale, de forma trazable y a escala, la
información de capacidad sobre la que analizan. La segunda es que la
interacción entre límites de uso y \textit{rates} queda definida como
semántica del modelo (la intersección de rectas que da lugar al
BoundedRate), pero no caracterizada sobre ningún dominio real. Qué formas
adopta esa relación en las APIs de un sector concreto, y qué
consecuencias tiene para el consumidor, es una pregunta que ninguno de
estos trabajos aborda.

\section{El canal API del ecosistema: ACM, Datasheets y motor de
análisis}
\label{sec:estado:acm}

Esta sección revisa los dos Trabajos Fin de Máster presentados al
principio del capítulo, que extienden el ecosistema de Pricing
Intelligence al canal API. Son el de Gavira Sánchez \citep{gavira2026tfm}
y el de Ruiz López, centrado en el agente HARVEY4API \citep{ruizlopez2026tfm}.

El \textbf{ACM} (\textit{API Capacity Model}) reconceptualiza la
capacidad para superar las tres limitaciones señaladas al final de la
Sección~\ref{sec:estado:capacidad}. Frente a la visión unidimensional,
el ACM define la capacidad como un conjunto de \textit{dimensiones}
independientes, cada una expresada sobre su propia métrica:
\textit{Volume} (las cuotas), \textit{Frequency} (los \textit{rates}),
\textit{Concurrent Requests} y \textit{Results Per Request}. A ellas se
añade el CRF (\textit{Consumption Request Factor}), que no es un límite
sino la regla que declara cuántas unidades de cada métrica gasta una
petición. En una API de correo esa distinción importa. Una misma llamada
consume una petición y, a la vez, tantos correos de la cuota mensual como
destinatarios incluya \citep{gavira2026tfm}. Frente a la semántica
delegada en el limitador, el ACM decide la clasificación comparando los
límites entre sí. De los que una API publica, el de menor unidad temporal
es el \textit{rate}, que gobierna la velocidad, y el de mayor unidad es
la cuota, que gobierna el volumen. Los 900 puntos por minuto y las 5.000
peticiones por hora de la API de GitHub \citep{githubratelimits} quedan
así clasificados sin necesidad de conocer el mecanismo interno del
proveedor \citep{gavira2026tfm}. La interpretación
geométrica de la Sección~\ref{sec:estado:analisis-capacidad} queda
integrada en el modelo como semántica oficial de la capacidad.

La serialización del modelo, \textbf{ACM2YAML}, da lugar a la
\textbf{Datasheet}: un documento orientado a máquinas que reúne en una
sola representación la capacidad de una API. En el canal API esa
capacidad aparece dispersa entre páginas de precios, documentación de
desarrollador, especificaciones OpenAPI y políticas de uso. Su
estructura combina los conceptos de capacidad con los elementos de
pricing que los contextualizan. Los primeros son los catálogos de cuotas
y rates (bloques \texttt{capacity} y \texttt{max\_power}), el CRF
serializado como \texttt{workload} y el sobrecoste como
\texttt{overage\_cost}. Los segundos, la jerarquía
Plan~$\rightarrow$~Endpoint~$\rightarrow$~Alias. Frente al plan como
único eje de segmentación, esa jerarquía asocia cada restricción al
nivel al que aplica, sin crear planes artificiales \citep{gavira2026tfm}.
La versión operativa, ACM2YAML v0.3, captura dos de las cuatro
dimensiones del modelo (Volume y Frequency, además del CRF) y deja fuera
concurrencia y resultados por petición. Su propio autor la califica de
versión operativa parcial, priorizada para poner en marcha el ecosistema
antes que para agotar la expresividad del metamodelo
\citep{gavira2026tfm}.

Sobre la Datasheet operan las herramientas del canal.
\textbf{PRIME4API} refactoriza la biblioteca Pricing4API como API REST y
convierte el análisis en un proceso dirigido por la Datasheet: cada
pregunta se reduce a operaciones sobre el BoundedRate y la curva de
capacidad \citep{gavira2026tfm}.
Junto al motor, Gavira Sánchez define un catálogo de diez preguntas de
análisis (Q1--Q10) en tres categorías: análisis de la capacidad de una suscripción, selección de la suscripción
óptima y equilibrio entre cuotas heterogéneas. Lo instancia sobre tres
pricings de dos servicios (SendGrid, en dos versiones, y DBLP), y esas
respuestas de referencia funcionan como \textit{ground truth} para validar el
ecosistema \citep{gavira2026tfm}. La tercera categoría es especialmente
relevante para este trabajo. Q8 (cuál de dos cuotas se agota primero
dado un CRF), Q9 (qué CRF iguala el agotamiento de dos cuotas) y Q10
(cuánto cuesta completar un volumen incurriendo en sobrecoste) razonan
sobre la interacción entre límites heterogéneos. Esa categoría no tiene correspondencia en el
catálogo de Fresno-Aranda precisamente porque su modelo carecía del
factor de conversión \citep{gavira2026tfm}. Por último, HARVEY4API expone
estas capacidades tras una interfaz conversacional, publicando las
operaciones del motor como herramientas mediante el \textit{Model
Context Protocol} \citep{mcp,ruizlopez2026tfm}. La extensión de SPHERE
al canal API aporta colecciones de Datasheets, renderizado en tiempo
real y edición validada \citep{gavira2026tfm,sphere}. La extracción
automática es tarea de otra herramienta del canal, A-MINT4API, y el
editor de SPHERE sirvió para validar visualmente las Datasheets que
genera \citep{gavira2026tfm}.

Ruiz López ha desarrollado también dos \textit{skills} para caracterizar cómo se accede a un servicio SaaS \citep{ruizlopez2026skills}.
Una \textit{skill} es un conjunto de instrucciones empaquetadas que un modelo de lenguaje sigue para una tarea.
La primera, \texttt{provider-service-resolver} (versión 1.1), recibe el nombre de un SaaS e identifica a su proveedor.
Después clasifica cada candidato de su catálogo como servicio autónomo, servicio dependiente, complemento, característica u oferta, y cada afirmación lleva su fuente oficial con el extracto literal.
La segunda, \texttt{saas-interface-channel-analyzer} (versión 3.0), parte de ese resultado y separa la interfaz del canal.
Una interfaz es un mecanismo definido por el proveedor con un solo paradigma de interacción, y el canal es ese paradigma: programático, por eventos, interactivo, conversacional o agéntico.
Con esas reglas, el envío por SMTP y la API REST de un proveedor de correo son dos interfaces del mismo canal programático, porque cambian el protocolo y la autenticación pero no el paradigma.

Esta distinción es un avance reciente del trabajo.
Hasta octubre de 2026, esta memoria llamaba canal a cada vía de acceso a un servicio: el canal web, donde se publica el pricing, y el canal API, que es el que se modela.
Esa palabra mezclaba dos cosas, la vía de acceso y la forma de interactuar a través de ella.
En la revisión con el tutor del 15 de octubre de 2026 se acordó usar interfaz con la definición del NIST: la frontera común entre sistemas o módulos independientes en la que tienen lugar las interacciones \citep{nist80053r5}.
Con ella, la API REST, el relé SMTP y el panel web de un proveedor son interfaces distintas del mismo servicio.
El canal queda para el paradigma de interacción, como en las \textit{skills} de Ruiz López.
La multicanalidad de un servicio de correo se lee entonces como el conjunto de sus interfaces y el canal de cada una.
Las expresiones canal web y canal API que siguen en esta memoria nombran, por tanto, la interfaz web y las interfaces programáticas del servicio.

El resultado es una cadena completa y operativa: extracción, modelo,
motor, agente y hub funcionan sobre ACM2YAML v0.3. Sus límites, sin
embargo, están declarados con la misma claridad en la propia memoria de
Gavira Sánchez \citep{gavira2026tfm}. La validación descansa sobre tres
pricings de dos servicios. Dos de las cuatro dimensiones de capacidad carecen de
serialización. La integración entre capacidad y estructura comercial se
resolvió de forma pragmática y se señala como trabajo futuro. Y tanto
completar la serialización como someterla a validación formal figuran
literalmente entre sus líneas pendientes. Con la plataforma construida,
falta confrontarla con un dominio completo.

\section{El dominio de aplicación: las APIs de correo electrónico}
\label{sec:estado:correo}

El dominio elegido para esa confrontación es el del correo electrónico
programático, presentado en la
Sección~\ref{subsec:introduccion:canales}. Es un dominio maduro y
competitivo, en el que varios proveedores ofrecen el mismo servicio por
API. La información de
capacidad (\textit{rates} por endpoint, cuotas por plan, factores de
consumo por destinatario o por lote) resulta especialmente dispersa y
heterogénea. Eso lo convierte en un banco de pruebas exigente para el
instrumental recién revisado.

Sobre cómo se comparan estas APIs se han revisado dos comparativas
públicas del sector. Las dos las publican actores que ofrecen servicios
de correo y que se incluyen a sí mismos en la comparación
\citep{mailtrap2026comparativa,mailslurp2026comparativa}. Adoptan el formato de listas de
recomendaciones: útiles como panorámica del mercado, pero elaboradas con
criterios no declarados, sin trazabilidad de fuentes, no reproducibles y
no procesables por máquina. En la literatura académica revisada en este
capítulo, las APIs de correo aparecen únicamente como casos de
ilustración o validación. SendGrid, en dos versiones, forma parte del
\textit{ground truth} sobre el que se validó el canal API del
ecosistema, y MailerSend sirvió como banco de pruebas en el diseño de
ACM2YAML \citep{gavira2026tfm}. En los trabajos revisados en este
capítulo no se ha identificado un tratamiento sistemático del dominio
como objeto de comparación. Tampoco ofrecen un
catálogo estructurado, trazable y procesable de las condiciones de
consumo de las APIs de correo, ni un método asentado para compararlas
combinando requisitos funcionales y no funcionales.

\section{Análisis del estado del arte}
\label{sec:estado:analisis}

La Tabla~\ref{tab:estado:sintesis} sintetiza la revisión anterior. Para
cada propuesta se indica el canal que modela, cómo entiende la
capacidad, qué análisis ofrece y sobre qué se ha validado.

\begin{table}[H]
\centering
\small
\begin{tabular}{|p{3.1cm}|p{1.9cm}|p{2.9cm}|p{2.9cm}|p{2.6cm}|}
\hline
\textbf{Propuesta} & \textbf{Canal} & \textbf{Capacidad} &
\textbf{Análisis} & \textbf{Validación} \\
\hline
iPricing / Pricing2YAML
\citep{garciafernandez2024pricing4saas,garciafernandez2024racing} &
Web & Límites de uso sobre \textit{features} & 7 operaciones sobre
CSOP \citep{garciafernandez2025ipricing} & $>$150 pricings; corpus
SPECTRA \citep{spectra2025} \\
\hline
Pricing4APIs / SLA4OAI
\citep{fresno2025pricing4apis,gamezdiaz2019sla4oai} & API &
Unidimensional (peticiones); cuota/\textit{rate} según ventana del
limitador & Validación de consistencia (VC1--VC4) & Casos de
ilustración \\
\hline
Smart LAMA \citep{fresno2022lama,fresno2023gecon} & API (arquitectura) &
Máximo \textit{workload} de la LAMA & CSOP: max.\ peticiones, min.\
coste, min.\ tiempo & Notebook con 10 preguntas \\
\hline
Catálogo de operaciones \citep{fresno2023catalogo} & API & La del modelo
Pricing4APIs & 11 operaciones formuladas, sin implementación & --- \\
\hline
Pricing4API (biblioteca) \citep{pricing4api,gavira2025tfg} & API &
Curva de capacidad $C(t)$; BoundedRate & Q1--Q7, programático y visual &
Casos sin entrada estructurada \\
\hline
ACM / Datasheet / PRIME4API \citep{gavira2026tfm,ruizlopez2026tfm} & API
& 4 dimensiones + CRF (v0.3 serializa 2) & Q1--Q10 dirigido por la
Datasheet & 3 APIs (\textit{ground truth}) \\
\hline
\end{tabular}
\caption{Síntesis del estado del arte revisado.}
\label{tab:estado:sintesis}
\end{table}

De esta síntesis se desprenden cuatro observaciones, que son los cuatro
huecos sobre los que se articula este trabajo.

\textbf{Primera: el instrumental existe; su validación a escala de
dominio, no.} La cadena completa del canal API funciona, pero se ha
ejercitado sobre tres pricings de dos servicios \citep{gavira2026tfm}. Ningún
trabajo ha confrontado todavía el modelo con un conjunto representativo de las APIs de un sector, elegido con un criterio declarado,
incluidas las que no encajan con
naturalidad. Con una confrontación así se contrasta si la Datasheet es un formato
general o un formato que funciona donde ya se probó. Este es el
hueco que aborda el catálogo de Datasheets de correo (contribución
C-01).

\textbf{Segunda: la expresividad de ACM2YAML v0.3 es parcial, y sus
carencias están declaradas pero no inventariadas.} Se sabe qué
dimensiones faltan (concurrencia, resultados por petición), pero no qué
extensiones concretas exige un dominio real ni con qué prioridad. Ese
inventario solo puede producirlo la confrontación sistemática con casos
reales, documentando cada carencia con la API que la revela. Es el
hueco de la evolución justificada del esquema (C-02).

\textbf{Tercera: la relación entre límites de uso y \textit{rates} está
definida, pero no caracterizada.} La semántica existe (la intersección
que define el BoundedRate) y las operaciones Q8--Q10 la ejercitan sobre
casos concretos. Lo que ningún trabajo ofrece es una caracterización de
esa relación sobre un dominio real: qué patrones adopta, cuándo el
\textit{rate} enmascara a la cuota o la cuota al \textit{rate}, y qué
implicaciones tiene para el consumidor. Es el hueco de la contribución
C-03.

\textbf{Cuarta: no hay método de comparación entre APIs.} Los catálogos
de preguntas existentes (las once operaciones de Fresno-Aranda, las diez
de Gavira Sánchez) razonan sobre la capacidad de un pricing. La decisión
que motiva este trabajo, \textit{qué API conviene a unas necesidades
dadas}, exige además preguntas funcionales y no funcionales propias
del dominio y un procedimiento sistemático para responderlas sobre un
conjunto de Datasheets comparables. Ese método, con su catálogo de
preguntas de decisión, es la contribución C-04.

Este trabajo, conviene subrayarlo, no propone un modelo nuevo ni un
motor nuevo. Aplica y extiende los instrumentos revisados en este
capítulo (la Datasheet, PRIME4API, SPHERE), llevándolos a un dominio
real y a escala. El catálogo resultante queda en el repositorio de datos
que acompaña al trabajo (Sección~\ref{sec:conclusiones:replicacion}),
preparado como contribución a SPHERE, y el Capítulo~\ref{cap:evolucion}
propone las extensiones del esquema que el dominio justifica. Los capítulos
siguientes desarrollan cómo.

Este trabajo aporta a SPHERE dos piezas.
La primera es una colección con las ocho APIs del catálogo como \textit{API products}, es decir, interfaces programáticas con un pricing propio.
Cada una lleva su Datasheet y su pricing en Pricing2YAML \citep{garciafernandez2025trends}.
La segunda es una propuesta de sección de términos de consumo para la ficha de cada API.
Los términos de consumo (\textit{API consumption terms} \citep{fisapuntes}) son los términos y condiciones del proveedor que afectan a la velocidad y a la capacidad de consumo de su API, como sus \textit{rates}, sus cuotas y su política de uso.
La sección llevará una pestaña con las versiones registradas de esos términos y la procedencia de cada dato: URL, tipo de fuente y fecha.
El repositorio del trabajo incluye una maqueta de la colección con esa pestaña y con las versiones anteriores de cada pricing, tomadas sobre todo de capturas de Wayback Machine.
SPHERE y Pricing2YAML son del grupo ISA: la colección los usará tal como son, y la sección se propone como extensión de la plataforma.

Las \textit{skills} de Ruiz López se aplican sin modificarlas, con el procedimiento de la Sección~\ref{sec:propuesta:criterios}.
